% Owner: Hassan Ashraf 24L3073
\section{Introduction}
\label{sec:intro}

ProcurePulse is an AI-assisted tool that helps businesses decide \emph{when} to
buy metals and whether a supplier's price is fair. It studies the prices of
gold and other metals traded on COMEX, predicts the short-term direction of
gold, and compares a supplier's quote with the expected price range to
recommend \textbf{Buy}, \textbf{Hold} or \textbf{Negotiate}. The project also
tests whether external economic signals, such as the US dollar, interest
rates, oil and inflation, help to predict gold.

This is the final report of the project. It brings together the three
deliverables, whose objectives were to:
\begin{enumerate}
  \item define the real-world problem, the target users and why machine
        learning is needed (Phase~1);
  \item collect 16 years of daily market data from public sources with
        reproducible, logged Python scripts, then clean it, engineer features
        and explore it visually (Phase~2);
  \item build, tune and compare classical machine learning models and a deep
        learning model for the five-day direction of gold (Phase~3);
  \item evaluate the models with robust validation and suitable metrics, and
        analyse their errors and bias (Phase~4);
  \item turn the results into Buy, Hold and Negotiate recommendations and
        deliver them through an application with real-time inference
        (Phase~5, Track~A);
  \item document the whole project in one report (Phase~6).
\end{enumerate}

The report is organised as follows. Section~\ref{sec:problem} defines the
problem, the target audience and the AI justification. Section~\ref{sec:collection}
describes the data sources and the raw data, Section~\ref{sec:cleaning}
explains how the data was cleaned, and Section~\ref{sec:features_eda}
presents the features and the exploratory data analysis.
Sections~\ref{sec:classical} and~\ref{sec:deep} describe the classical models
and the LSTM, Section~\ref{sec:evaluation} compares all models on the test
years, and Section~\ref{sec:error-bias} analyses their errors and bias.
Section~\ref{sec:recommendation} explains how quotes are turned into
recommendations, and Section~\ref{sec:application} presents the application.
Section~\ref{sec:contributions} shows the system architecture and who built
each part, and Section~\ref{sec:conclusion} summarises the results and future
work. All code is in the project's GitHub repository.
